%---------------------------------------导包------------------------------------------
\documentclass[12pt,a4paper,AutoFakeBold]{article} %字号为小四，纸张大小为a4，文档类型为article学术类，AutoFakeBold开启伪粗体功能
\usepackage[margin=1in]{geometry} %设置边距，符合Word设定
%\usepackage[left=2.50cm,right=2.50cm,top=2.80cm,bottom=2.50cm]{geometry} %页边距设置

%%字体与行距
\usepackage{setspace} %修改行距
\usepackage{fontspec} %字体设置
\usepackage{xeCJK} %中文包
\usepackage{indentfirst} %中文首行缩进
\usepackage[UTF8,heading=true]{ctex} %修改标题样式

%%图表、公式包
\usepackage{graphicx} %插入图片
\graphicspath{{Figures/}} %文章所用图片在当前目录下的Figures目录
%\usepackage{subfigure} %并排子图，subcaption包含了subfigure
\usepackage{caption}
\usepackage{subcaption} %这两个包是为了调整图表标题和子标题
\usepackage{booktabs} %画三线表
\usepackage{amsmath,amsthm,amssymb,amsfonts} %数学公式包

%%摘要，目录，页眉页脚
\usepackage{abstract} %调整摘要格式
\usepackage{lipsum} %随机从某著作中挑选段落生成，\lipsum[2]
\usepackage{titlesec}
\usepackage{titletoc} %这两个包是为了修改目录样式
\usepackage{fancyhdr} %设置全文页眉、页脚的格式

%%其他
\setlength{\headheight}{27pt} %解决警告中headheight过小的问题 

%--------------------------------------前期设置-----------------------------------------
%设置文中的各种引用和链接

%%目录
\usepackage{hyperref} %对目录生成链接，注：该宏包可能与其他宏包冲突，故放在所有引用的宏包之后
\usepackage{cleveref} %实现图片和表格、公式的引用，cleveref包必须放在hyperref包之后
\renewcommand{\contentsname}{\centerline{\heiti\zihao{-3} 目录}} %经过设置word格式后，居中，黑体，小三，加粗

\titlecontents{section}[0cm]{\bfseries\songti\fontsize{14pt} %1级标题，宋体四号，加粗，标题与页面距离，下面的解释可能有误
{\baselineskip}\selectfont}{\hspace*{1em}\contentslabel{1em}\ } %\hspace*{1em}调整标号与页面距离，\contentslabel{1em}调整标号与标题之间间距，
{}{\titlerule*[0.5pc]{$\cdot$}\contentspage\hspace*{0cm}} %\titlerule*[0.5pc]{$\cdot$}设置标题后引线，\contentspage\hspace*{0cm}页码末尾与页面距离

\titlecontents{subsection}[0cm]{\bfseries\songti\fontsize{12pt}	%1.1级标题，宋体小四
{\baselineskip}\selectfont}{\hspace*{3em}\contentslabel{1.5em}\ }
{}{\titlerule*[0.5pc]{$\cdot$}\contentspage\hspace*{0cm}} 

\titlecontents{subsubsection}[1cm]{\bfseries\songti\fontsize{10pt} %1.1.1级标题，宋体五号
{\baselineskip}\selectfont}{\hspace*{3em}\contentslabel{2.5em}\ }
{}{\titlerule*[0.5pc]{$\cdot$}\contentspage\hspace*{0cm}}

%%链接和PDF属性
\hypersetup{colorlinks = true,  %将链接文字带颜色
	linkcolor = black, %图表引用颜色设置为蓝色
	anchorcolor = black,
    citecolor = blue, %文献引用颜色设置为蓝色
	bookmarksopen = true, %展开书签
	bookmarksnumbered = true, %书签带章节编号
	pdftitle = Learning Notes, %标题
	pdfauthor = Xiao} %作者

%%参考文献
\bibliographystyle{plain} %参考文献引用格式
\newcommand{\upcite}[1]{\textsuperscript{\cite{#1}}} %引用参考文献时，\upcite{1}，能标在右上角

%%大标题
\title{\songti\zihao{3}\textbf{20221023\textasciitilde20221029学习笔记}} %\textasciitilde表示波浪号
\author{}
\date{}

%%章节标题
\ctexset{ 
	section={
		format=\bfseries\songti\zihao{-3} %设置section标题为宋体、小三号字，\bfseries表示加粗
	},
	subsection={
		format=\bfseries\songti\zihao{4} %设置subsection标题为宋体、4号字
	},
	subsubsection={
		format=\bfseries\songti\zihao{-4} %设置subsubsection标题为宋体、小四号字
	}
}

%%摘要
%\renewcommand{\abstracttextfont}{\fangsong} %摘要内容字体为仿宋
\renewcommand{\abstractname}{\bfseries\songti\zihao{4} 摘 \quad 要} %更改摘要二字的样式，四号，加粗

%%全文字体与行距
\setmainfont{Times New Roman} %数字英文字体
\renewcommand{\baselinestretch}{1.5} %设置为1.5倍行距

%%页眉页脚
\pagestyle{fancy} %让页眉自动显示节标题

%%图表标注
\captionsetup[figure]{font=small,labelsep=quad} %space去掉点,period加点,不加space、period这两个就是冒号,quad是空格,small对应5号字体
\captionsetup[table]{labelsep=quad} %使用\captionsetup[figure]{name={Fig.},labelsep=period}图片标题会变成：Fig.XXXX
%\captionsetup[table]{font={small,bf,stretch=1.25}, justification=raggedright} %small对应5号字体，bf表示加粗，stretch=1.25设置行距为1.25，justification=raggedright表示左对齐

%--------------------------------------文章开始-----------------------------------------
\begin{document}
	
	\thispagestyle{empty} %设置本页无页眉页脚页码

	%摘要
	% \begin{abstract}
	% 	\lipsum[2]

	% 	\noindent{【{\bfseries\songti 关键词}】啦啦 \quad lala} %\noindent表示取消缩进
	% \end{abstract}
	
	%目录
	\tableofcontents
	
	%正文开始----------------------------------------------------------------------------
	\maketitle %生成文章的大标题
	\section{学习进度}

		这周。

	\newpage
	\section{区别与优缺点}
		\subsection{区别}

			{\bfseries\songti 啦啦啦}结构示意图如图\ref{pho:structure}所示。
		
			\begin{figure}[htbp]
				\centering
				\includegraphics[scale=0.6]{structure.png}
				\caption{结构对比}
				\label{pho:structure}
			\end{figure}

			\subsubsection{区别区别}

				这句话是测试能否进行引用及支持中文\upcite{1}。

				可由公式\eqref{equ:q1-1-Loss}表示。
				
				对$Loss$函数求偏导得到线性方程组。
				
				通过矩阵法可求解常数项和系数(即$b_{0},b_{1},b_{2},...b_{n}$)。

				\begin{equation}
					Loss = \sum_{i=1}^{n} (y_i-\widehat{y_i})^2 
					\label{equ:q1-1-Loss}
				\end{equation}

				求解得到系数如表\ref{tab:q1-humidity coefficient}所示：
				\begin{table}[!htbp]
					\caption{系数}%表格标题
					\centering         
					\begin{tabular}{cccccccc}
						\toprule[1.5pt]%表格顶部的粗线
						深度 & $\beta _{0}$  &  $\beta _{1}$  &   $\beta _{2}$  &   $f_{0}$  &   $f_{1}$&   $f_{2}$&   $f_{3}$\\
						\midrule[1pt]%表格中部的粗线
						10cm&1&2&3&4&5&6&7\\
						20cm&1&2&3&4&5&6&7\\
						40cm&1&2&3&4&5&6&7\\
						\bottomrule[1.5pt]%表格底部的粗线
					\end{tabular}
					\label{tab:q1-humidity coefficient}%标签得放在tabular里面或后面才能正常引用
				\end{table}

	%参考文献
	\newpage
	\addcontentsline{toc}{section}{参考文献} %将参考文献加入目录
	\bibliography{books}

\end{document}